\chapter*{\texorpdfstring{Guia de lectura i itineraris\\ d'aprenentatge}{Guia de lectura i itineraris d'aprenentatge}}
\addcontentsline{toc}{chapter}{Guia de lectura}
\markboth{Guia de lectura}{Guia de lectura}

\begingroup
\renewcommand{\thesection}{\arabic{section}}
\setcounter{section}{0}
\setlength{\emergencystretch}{1.5em}

Aquest capítol és un mapa per a la primera lectura. No cal aprendre
ara els noms de totes les construccions que hi apareixen: les trobarem
definides al seu lloc. La seva funció és mostrar què aporta cada etapa,
com es relacionen les quatre parts del llibre i com podem comprovar
que l'estudi ens està fent avançar.

La introducció explica \emph{per què} val la pena aprendre categories.
Aquí ens ocuparem de \emph{com} fer-ho. Podeu llegir aquesta guia
abans de començar i tornar-hi en acabar un bloc de capítols, quan
els noms del mapa ja s'hagin convertit en construccions familiars.


\section{El punt de partida i la destinació}

El llibre no pressuposa coneixements previs de teoria de categories.
Sí que pressuposa una familiaritat elemental amb el llenguatge
matemàtic: conjunts, aplicacions, composició, relacions i
demostracions senzilles. Per començar, convé poder fer quatre coses:
\begin{itemize}
  \item distingir el domini i el codomini d'una aplicació i calcular
        una composició;
  \item treballar amb productes cartesians, subconjunts i antiimatges;
  \item llegir quantificadors i separar una afirmació d'existència
        d'una afirmació d'unicitat;
  \item distingir una implicació, el seu recíproc i una equivalència,
        i entendre el paper d'un contraexemple.
\end{itemize}
Si alguna d'aquestes operacions és poc familiar, és millor dedicar-hi
una breu revisió que atribuir-ne la dificultat a l'abstracció categòrica.

No cal dominar prèviament tots els àmbits que proporcionen exemples.
Els grups, els espais vectorials i la topologia enriqueixen la lectura,
però, si encara no els coneixeu, podeu començar amb conjunts, preordres
i grafs. Els exemples del dual i del bidual demanen àlgebra lineal;
els que utilitzen continuïtat demanen nocions de topologia. Es poden
reprendre en una segona lectura sense abandonar el fil general.
La lògica aporta exemples des del començament, però no cal conèixer
ja la lògica categòrica per seguir-los.

La destinació és una comprensió inicial de les categories, els
functors i la naturalitat, seguida de les propietats universals,
Yoneda, els límits i les adjuncions. Aquest bagatge permet arribar
al tancament cartesià, als subobjectes i a una primera introducció
als topos. El volum prepara l'estudi posterior de la lògica categòrica;
no pretén desenvolupar-ne aquí una teoria completa.

\section{Quatre parts per a un mateix aprenentatge}

El llibre separa físicament la \textbf{Teoria}, la \textbf{Pràctica},
els \textbf{Exercicis proposats} i els \textbf{\NomTests}.
Aquesta separació facilita la consulta, però no obliga a estudiar
primer tota la teoria i només després tota la pràctica. Per a una
primera lectura, és més útil treballar \emph{per temes}: llegir un
capítol teòric, estudiar els problemes resolts corresponents,
intentar exercicis i comprovar la comprensió amb el test del mateix
tema.

\begin{center}
\renewcommand{\arraystretch}{1.2}
\begin{tabularx}{\linewidth}{@{}>{\raggedright\arraybackslash}p{0.25\linewidth}
                                  >{\raggedright\arraybackslash}X@{}}
\toprule
\textbf{Part} & \textbf{La pregunta que ajuda a respondre} \\
\midrule
Teoria & Quina és la definició, per què té sentit i què se'n pot deduir? \\
Pràctica & Com s'utilitza en un problema i com s'escriu l'argument? \\
Exercicis proposats & Puc reconèixer i aplicar la idea sense una solució completa? \\
\NomTests & Distingeixo el concepte de les confusions que s'hi assemblen? \\
\bottomrule
\end{tabularx}
\end{center}

Els nou primers temes tenen capítols corresponents en les quatre
parts. La numeració torna a començar en cada part; per això, quan
cerqueu material d'un tema, fixeu-vos tant en el número com en la
part i el títol. El capítol teòric desè, \emph{Pont cap a la lògica
categòrica}, és un epíleg d'orientació i no té un capítol paral·lel
de pràctica, exercicis o tests.

Els blocs d'obertura dels capítols teòrics indiquen prerequisits
i objectius. Llegiu-los com una invitació a activar coneixements
anteriors. Els blocs de síntesi del final recullen les idees centrals,
els errors típics i lectures recomanades. Serveixen per revisar,
però no substitueixen els exemples i les demostracions del capítol.

\section{El mapa conceptual dels capítols}
\label{sec:mapa-lectura}

\ifpublicacioparcial
\noindent\textbf{Nota sobre aquesta edició.} Aquesta versió parcial
només inclou els capítols 1--3. El mapa descriu, tanmateix, els deu
capítols del volum complet, perquè pugueu situar el que estudieu
dins del recorregut sencer.
\fi

El recorregut es pot agrupar en quatre etapes. Els capítols 1--3
construeixen el llenguatge i els nivells de comparació; els 4--6
organitzen les construccions mitjançant propietats universals;
els 7--9 combinen aquestes eines en estructures amb lectura lògica;
el 10 mostra la continuació. Aquesta divisió és una ajuda per orientar-se,
no una frontera rígida: les idees de les primeres etapes reapareixen
constantment en les següents.

\subsection*{1. Categories: aprendre la gramàtica de les fletxes}

El primer capítol introdueix objectes, morfismes, identitats i
composició. També ensenya a llegir diagrames, construeix exemples
i diferencia isomorfismes, monomorfismes, epimorfismes, seccions
i retraccions. La categoria oposada i la dualitat permeten
reutilitzar definicions i arguments en l'altra direcció.

És un capítol extens perquè fixa els hàbits que necessitarem després.
No convé voler memoritzar tots els exemples alhora. Escolliu-ne uns
quants i pregunteu sempre què són els objectes, què són les fletxes
i com es componen. Les construccions de categories noves i les
precaucions de mida amplien aquest vocabulari.

\noindent\textbf{Comprovació de progrés.} Podeu verificar els axiomes
en un exemple concret i explicar per què una fletxa és, o no és,
un isomorfisme? Podeu escriure l'equació d'un quadrat commutatiu
sense invertir l'ordre de la composició?

\subsection*{2. Functors: passar d'un context a un altre}

El segon capítol puja un nivell: no compara només objectes d'una
categoria, sinó categories entre elles. Els functors actuen sobre
objectes i morfismes i preserven identitats i composició. El graf
subjacent i la categoria lliure permeten veure una parella de
construccions que reapareixerà com a adjunció.

Cal dedicar una atenció especial als homfunctors. $\Hom(A,-)$
actua per postcomposició; $\Hom(-,B)$, per precomposició i amb
inversió de sentit. Aquesta variància no és un detall de notació:
és la base de la representabilitat i de Yoneda. Les nocions de
functor ple i fidel descriuen l'acció sobre cada homconjunt.

\noindent\textbf{Comprovació de progrés.} Per definir un functor,
podeu donar-ne l'acció sobre les fletxes, no només sobre els objectes?
Podeu explicar què fa $\Hom(-,B)$ amb una fletxa $X'\to X$?

\subsection*{3. Transformacions naturals i equivalències: exigir coherència}

El tercer capítol compara functors. La naturalitat exigeix que
una família de morfismes sigui compatible amb totes les fletxes
del domini. Els exemples de la interpretació de fórmules i del
bidual mostren per què aquesta compatibilitat és matemàticament
significativa, no només una convenció gràfica.

Les categories de functors permeten tractar aquestes comparacions
com a morfismes. L'equivalència de categories precisa quan dos
contextos descriuen la mateixa estructura essencial sense tenir
la mateixa presentació. La composició horitzontal afegeix un
segon tipus de composició que cal distingir de la vertical.

\noindent\textbf{Comprovació de progrés.} Podeu dibuixar el quadrat
de naturalitat d'una família proposada i comprovar-lo per a una
fletxa arbitrària? Podeu distingir un isomorfisme d'objectes,
un isomorfisme natural i una equivalència de categories?

\subsection*{4. Propietats universals, representabilitat i Yoneda:
caracteritzar pel paper}

Aquí canvia el centre de gravetat. En lloc de construir un objecte
i estudiar-ne després les propietats, busquem un objecte que
resolgui una necessitat expressada amb fletxes. Els objectes
inicials i terminals ofereixen els primers casos; les categories
coma ajuden a reunir dades i compatibilitats en un mateix context.

La representabilitat expressa aquestes caracteritzacions amb
homfunctors. Yoneda fa precisa la idea de conèixer un objecte
pel sistema de morfismes que hi arriben. És normal que aquesta
etapa demani més d'una lectura: cal coordinar categories, functors,
transformacions i elements d'homconjunts.

\noindent\textbf{Comprovació de progrés.} Podeu demostrar que dos
objectes terminals són isomorfs i indicar on s'utilitza la unicitat?
En una representació, podeu distingir el functor representat,
l'objecte representant i la dada universal? Podeu seguir la
bijecció de Yoneda a partir de l'avaluació en la identitat?

\subsection*{5. Límits i colímits: reunir construccions en un patró}

Un diagrama fixa una configuració d'objectes i morfismes; un con
és una manera compatible de relacionar-s'hi. El límit és el con
universal. Segons la forma del diagrama, obtenim productes,
equalitzadors, pullbacks o l'objecte terminal. Els colímits en
són la versió dual.

El pullback mereix una atenció especial, perquè reapareixerà
en la reindexació i els subobjectes. També cal separar dues
preguntes: si una categoria té un límit i si un functor el
preserva. No són la mateixa afirmació.

\noindent\textbf{Comprovació de progrés.} Podeu calcular un pullback
de conjunts i comprovar-ne la propietat universal? Podeu
distingir el diagrama, el con i el seu vèrtex, i explicar
què vol dir que un functor preserva aquest límit?

\subsection*{6. Adjuncions: dues descripcions de la mateixa dada}

Les adjuncions relacionen functors mitjançant una bijecció
natural d'hom\-con\-junts. El capítol comença amb preordres,
on la correspondència es llegeix com una equivalència de
desigualtats. Després introdueix unitat, counitat i identitats
triangulars, i mostra com les adjuncions expliquen la preservació
de límits i colímits.

Per a una primera aproximació, escolliu una adjunció i seguiu-la
en les dues descripcions: per bijeccions d'homconjunts i per
unitat i counitat. Les construccions lliures i la currificació
connecten aquest capítol amb exemples ja coneguts.

\noindent\textbf{Comprovació de progrés.} Podeu dir quins morfismes
es corresponen en una adjunció concreta i construir-ne els
transposats? Podeu recordar, amb una justificació conceptual,
que els adjunts per la dreta preserven límits i els de
l'esquerra preserven colímits?

\subsection*{7. Categories cartesianes tancades: funcions dins del context}

Els productes i els exponencials permeten parlar d'aparellament,
avaluació i currificació dins d'una categoria. El capítol
caracteritza el tancament cartesià i n'explica la relació
amb el càlcul lambda simplement tipat: els tipus es llegeixen
com a objectes i els termes, en context, com a morfismes.

El conjunt de totes les aplicacions $B\to C$ ofereix l'exemple
més familiar d'exponencial. Els preordres permeten reconèixer
la implicació com a adjunt de la conjunció. Els contraexemples
recorden que no tota categoria té aquesta estructura.

\noindent\textbf{Comprovació de progrés.} Podeu passar d'una fletxa
$A\times B\to C$ a la seva currificació $A\to C^B$ i tornar-ne
amb l'avaluació? Podeu explicar per què tenir productes no
garanteix tenir exponencials?

\subsection*{8. Subobjectes, imatges i factoritzacions: de les parts als predicats}

Un subobjecte no és simplement un subconjunt escrit amb un
altre nom: és una classe de monomorfismes amb un mateix
codomini, identificats de manera compatible. La reindexació
per pullback generalitza l'antiimatge i ofereix una lectura
estructural de la substitució.

Les interseccions expressen conjuncions de predicats; les
imatges i la regularitat permeten construir un adjunt
existencial de la reindexació. Aquesta etapa obliga a
controlar les hipòtesis: la regularitat no garanteix per
si sola unions de subobjectes ni un quantificador universal.

\noindent\textbf{Comprovació de progrés.} Podeu calcular una
reindexació a $\cat{Set}$ i interpretar-la com a substitució?
Podeu descriure la imatge d'un predicat sota una projecció
i explicar la relació $\exists_f\dashv f^*$?

\subsection*{9. Classificadors de subobjectes i topos: retrobar la representabilitat}

El classificador generalitza la funció característica d'un
subconjunt. La correspondència entre subobjectes d'$A$ i
morfismes $A\to\Omega$ recupera la representabilitat del
capítol quart. El topos elemental combina límits finits,
exponencials i classificador de subobjectes.

El capítol presenta $\cat{Set}$ i les categories de prefeixos
com a exemples. L'objectiu és comprendre les definicions
i la connexió entre les peces, no dominar encara la
teoria de feixos o la lògica interna dels topos.

\noindent\textbf{Comprovació de progrés.} Podeu reconstruir un
subconjunt com l'antiimatge del valor vertader per la
seva funció característica? Podeu enunciar les tres
condicions d'un topos i explicar què aporta cadascuna?

\subsection*{10. Pont cap a la lògica categòrica: saber què queda preparat}

L'epíleg situa les eines adquirides en un programa més ampli:
contextos, substitució, predicats, quantificadors, teories
i models. El mapa de fragments lògics i l'exemple de la
categoria sintàctica orienten sobre què es prolonga al
volum següent.

Llegiu-lo com una mirada retrospectiva i una invitació a
continuar, no com una exigència de dominar de cop tots els
temes anunciats. La sintaxi general, la correcció, la
completesa i les condicions de coherència pertanyen a
l'estudi posterior.

\noindent\textbf{Comprovació de progrés.} Quan l'epíleg parla de
models com a functors o de quantificadors com a adjunts,
podeu identificar les construccions d'aquest volum que
donen sentit a aquestes expressions?

\section{Com treballar un capítol}

Una lectura activa alterna comprensió, reconstrucció i aplicació.
Podeu utilitzar el cicle següent, adaptant-ne el ritme a la
dificultat del tema.

\begin{enumerate}
  \item \textbf{Situar la pregunta.} Llegiu l'obertura i els
        objectius del capítol. Escriviu en una frase quin
        problema intenta resoldre i quines idees anteriors
        necessita.
  \item \textbf{Desmuntar la definició.} Separeu les dades
        que s'han de donar, les operacions que s'hi afegeixen
        i les propietats que han de satisfer. Anoteu el
        domini i el codomini de cada fletxa.
  \item \textbf{Ancorar-la en exemples.} Treballeu un cas
        amb conjunts i un altre amb preordres o grafs.
        Quan pertoqui, afegiu-hi la lectura lògica. Busqueu
        també un exemple que no compleixi la definició.
  \item \textbf{Reconstruir una demostració.} Tanqueu el
        text i intenteu recuperar un argument curt. Després
        compareu-lo amb el llibre: on s'usa cada hipòtesi?
        On apareix l'existència i on la unicitat?
  \item \textbf{Fer servir la pràctica amb pausa.} Llegiu
        l'enunciat d'un problema resolt i proveu de fer-ne
        el primer pas abans de mirar la solució. Un cop
        llegida, torneu a escriure l'argument sense copiar-lo.
  \item \textbf{Intentar un exercici nou.} Comenceu pels
        de consolidació. Si us encalleu, identifiqueu el
        punt exacte i consulteu la indicació, no només
        una altra solució que s'hi assembli.
  \item \textbf{Tancar el cicle.} Feu el test, reviseu els
        errors i rellegiu la síntesi. Escriviu què podeu
        fer ara que no podíeu fer abans i què ha de quedar
        anotat per a una segona lectura.
\end{enumerate}

No cal reconstruir immediatament totes les proves llargues.
Però convé entendre'n l'estratègia i ser capaç de reproduir
arguments curts que reapareixen: provar una naturalitat,
usar la unicitat d'una propietat universal o construir
una fletxa amb el domini i el codomini correctes.

\subsection*{Una pauta per llegir diagrames}

Al llarg del llibre, quan una propietat categòrica té una expressió
natural amb fletxes, la presentem alhora amb un diagrama i amb les
equacions de composició que aquest expressa. Els diagrames formen
part del llenguatge matemàtic, no són una il·lustració opcional.
Aprofiteu aquesta doble presentació: passar del dibuix a les equacions
i de les equacions al dibuix és una destresa que cal practicar.

Davant d'un diagrama, no comenceu per recordar-ne el nom.
Seguiu aquesta pauta:
\begin{enumerate}
  \item Identifiqueu els objectes i els extrems de cada fletxa.
  \item Distingiu les dades donades de la fletxa que s'ha de construir.
  \item Escriviu els camins que tenen el mateix origen i el mateix destí.
  \item Traduïu-ne la commutativitat en equacions de composició.
  \item Si hi ha una propietat universal, digueu sobre quines
        dades es quantifica i què es demana que sigui únic.
\end{enumerate}

Una fletxa discontínua en un diagrama universal no es construeix
perquè el dibuix ho suggereixi: l'existència es justifica amb
una hipòtesi o una propietat. Igualment, un quadrat commutatiu
no és automàticament un pullback. Cal comprovar-ne també la
propietat universal.

\section{Els exemples recurrents com a punts de suport}

Canviar d'exemple no ha de significar tornar a començar.
Les mateixes famílies permeten veure cares diferents d'una idea.

\begin{description}
  \item[Conjunts.] Permeten calcular explícitament amb elements,
        productes, antiimatges, funcions i imatges. Són un
        primer laboratori, però no autoritzen a usar elements
        dins d'una categoria arbitrària.
  \item[Preordres.] Simplifiquen els homconjunts: hi ha com a
        màxim una fletxa entre dos objectes. Això converteix
        propietats universals en condicions d'ordre i adjuncions
        en equivalències de desigualtats. La unicitat automàtica
        és una ajuda, però també pot amagar una dificultat
        que reapareix en categories no primes.
  \item[Grafs.] Fan visible la diferència entre tenir fletxes
        i disposar de composició. Cal distingir la categoria
        de grafs, on les fletxes són homomorfismes de grafs,
        de la categoria lliure d'un graf, on són camins.
        La construcció lliure connecta els primers capítols
        amb les propietats universals i les adjuncions.
  \item[Lògica.] Permet interpretar deducció, conjunció,
        substitució i quantificació en els contextos adequats.
        Cal separar sempre la categoria prima que registra
        deduïbilitat de la perspectiva que conserva les proves,
        i també dels subobjectes que representen predicats.
\end{description}

Un quadern de correspondències pot ajudar: per a cada nova
construcció, escriviu què significa en dues d'aquestes famílies,
quina propietat comparteixen i quines hipòtesis necessita
cada realització. Si no existeix en un dels exemples, anoteu-ho:
aquest límit també forma part de la comprensió.

\section{Itineraris possibles}

Els itineraris següents són propostes de treball, no permisos
per ignorar les dependències. L'ordre del llibre continua sent
el recorregut més segur per a una primera lectura completa.

\subsection*{Una primera entrada: els tres capítols inicials}

Treballeu categories, functors i transformacions naturals amb
els seus materials paral·lels de pràctica, exercicis i tests.
La fita és poder moure-us entre els tres nivells: objectes i
morfismes, categories i functors, functors i transformacions.
No cal dominar encara tot el desplegament de la composició
horitzontal, però sí distingir-la de la vertical.

\ifpublicacioparcial
Aquesta edició parcial conté precisament aquests tres temes en
les quatre parts, a més dels materials complementaris. No és una
introducció completa als límits, les adjuncions o els topos;
les remissions als capítols absents s'han d'entendre com a
indicacions per continuar amb el volum complet.
\fi

\subsection*{Un recorregut complet d'autoaprenentatge}

Seguiu els nou temes amb el cicle de treball per capítols i
llegiu després el pont final. Cada cop que arribeu a una
definició general, torneu a un exemple anterior. Després
de Yoneda, reviseu els homfunctors; després de les adjuncions,
reviseu la construcció lliure; després del classificador,
reviseu la representabilitat. Aquestes tornades no són
repeticions inútils: fan visible la unitat del recorregut.

\subsection*{Un curs semestral}

Per a un curs d'un semestre, una selecció natural és treballar els capítols 1--6 com a nucli obligatori i triar, segons l'orientació del curs, els capítols 7--9. La teoria s'ha de combinar cada setmana amb una selecció de problemes resolts i exercicis proposats; els tests poden servir d'autoavaluació, no de substitut d'una prova escrita. El capítol 10 funciona bé com a sessió final de síntesi i d'obertura cap a la lògica categòrica.

\subsection*{Una lectura per al professorat}

Per preparar docència, convé començar per aquesta guia i pels blocs d'objectius i d'errors típics de cada capítol. Les parts de Pràctica permeten seleccionar exemples per a classe, mentre que els Exercicis proposats ofereixen material de treball autònom. Els quatre fils recurrents ---conjunts, preordres, grafs i lògica--- es poden modular segons el perfil de l'alumnat, però és recomanable conservar-ne almenys dos en paral·lel per evitar que les definicions quedin lligades a una sola categoria concreta.

\subsection*{Una lectura orientada a la lògica}

Manteniu el recorregut general, però doneu més pes a la
deduïbilitat del capítol 1 i als exemples d'interpretació
i substitució dels capítols 2 i 3. Als capítols 4--6,
consolideu representabilitat, pullbacks i adjuncions;
després, treballeu amb especial cura el tancament cartesià,
els subobjectes i el classificador.

No convé saltar directament a l'epíleg: expressions com
\enquote{el quantificador és un adjunt} orienten, però només es
comprenen quan sabem entre quines categories o preordres
actua l'adjunció i què afirma la seva correspondència.
Les equivalències lògiques no substitueixen la comprovació
de les hipòtesis categòriques.

\subsection*{Una lectura orientada als tipus i a la computació}

Feu servir els grafs i les categories lliures per entendre
la composició. Després de consolidar functors i naturalitat,
poseu l'accent en propietats universals, productes,
adjuncions, exponencials i currificació. El capítol 7
reuneix aquestes eines en la semàntica de tipus i termes.

Els capítols 8 i 9 poden constituir una segona etapa
si el primer objectiu és el càlcul lambda simplement tipat.
Caldrà reprendre'ls per arribar als predicats, els
classificadors i la perspectiva lògica més àmplia.

\section{Convencions tipogràfiques i de lectura}

Els noms de categories s'escriuen amb tipografia matemàtica, com $\cat{Set}$, $\cat{Graph}$ o $\cat{Cat}$. Els termes definits i les idees que cal retenir apareixen destacats dins dels entorns corresponents; el color reforça la jerarquia visual, però l'etiqueta textual ---\enquote{Definició}, \enquote{Teorema}, \enquote{Exemple}, \enquote{Exercici}, etc.--- n'indica sempre la funció. Les fletxes discontínues dels diagrames assenyalen habitualment el morfisme que una propietat universal afirma que existeix; el text adjacent n'explicita el paper. Els símbols amb més d'un ús, especialment $f^{*}$, es desambigüen pel context i es recullen a la taula de notació.

\section{Exercicis i tests: què comproven i què no}

Els exercicis proposats utilitzen una estrella,
$\star$, per assenyalar consolidació, i dues,
$\star\star$, per assenyalar ampliació. Comenceu pels
primers, però no interpreteu les marques com una mesura
exacta del temps que necessitareu. Una dificultat breu
de variància pot resultar més important que un càlcul llarg.

Les indicacions són una ajuda per reprendre el treball.
Abans de consultar-les, escriviu les dades, la conclusió
i almenys un intent. Després d'utilitzar una indicació,
torneu a fer l'exercici sense mirar-la. L'objectiu no és
haver reconegut una solució, sinó poder construir un argument.

Els tests tenen una única resposta correcta per pregunta.
\ifimprimible
En aquesta edició per imprimir, les respostes i les
justificacions de cada test es reuneixen al final del llibre
(apèndix \enquote{Solucions dels tests}). L'edició de pantalla
ofereix, a més, correcció automàtica en els lectors de PDF
compatibles.
\else
La correcció automàtica necessita un lector de PDF compatible
amb formularis AcroForm i JavaScript. Si el lector no
l'admet, podeu respondre igualment les preguntes i
consultar les respostes i les justificacions de cada test,
reunides al final del llibre (apèndix \enquote{Solucions dels
tests}), també usables en paper.
\fi
Consulteu-les només després d'haver respost tot el test.

Una resposta encertada no demostra per si sola que sapigueu
resoldre un problema. Per aprofitar el test, justifiqueu
per què l'opció triada és correcta i per què les altres
fallen. Si heu dubtat entre dues opcions, convertiu el
dubte en una pregunta per revisar a la teoria o a la
pràctica. Aquesta explicació és més útil que la puntuació
obtinguda sense justificació.

\section{Algunes dificultats que convé reconèixer aviat}

Hi ha errors que travessen diversos capítols. Identificar-los
permet corregir la manera de treballar, no només una resposta.

\begin{description}
  \item[Compondre sense comprovar els tipus.]
        Abans d'escriure $g\comp f$, verifiqueu que el
        codomini de $f$ és el domini de $g$. L'associativitat
        no fa componibles fletxes amb extrems incompatibles.
  \item[Confondre un exemple amb la definició general.]
        A $\cat{Set}$, mono i epi corresponen a injectivitat
        i exhaustivitat. En altres categories, les definicions
        són condicions de cancel·lació. Un mono que també
        és epi no és sempre un isomorfisme.
  \item[Tractar la naturalitat com una etiqueta.]
        Una família de fletxes no és natural només perquè
        s'hagi definit amb la mateixa fórmula. Cal escriure
        el quadrat i comprovar-lo per a tots els morfismes.
  \item[Oblidar les dades universals.]
        El producte inclou les projeccions; el límit inclou
        el con; l'exponencial inclou l'avaluació. La unicitat
        de l'isomorfisme es refereix a les compatibilitats
        amb aquestes dades, no a l'objecte nu.
  \item[Canviar l'orientació sense controlar-la.]
        La contravariància, la dualitat i la reindexació
        inverteixen direccions de maneres que cal precisar.
        Escriviu dominis i codominis abans de fer servir
        una regla recordada de memòria.
  \item[Afegir estructura que no s'ha suposat.]
        Tenir productes no implica tenir exponencials;
        la regularitat no dona per si sola disjuncions
        ni quantificació universal; un topos no és
        necessàriament booleà. Rellegiu les hipòtesis.
\end{description}

Si una demostració sembla bloquejada, pregunteu-vos primer
si el problema és de definició, de tipus, de variància o
d'hipòtesis. Moltes dificultats aparentment abstractes
es resolen quan precisem una d'aquestes quatre coses.

\section{Quan consultar els apèndixs i els materials finals}

Els apèndixs no constitueixen una etapa que s'hagi de
completar obligatòriament abans de començar la teoria.
Es poden consultar quan la lectura ho demani.

\begin{description}
  \item[Apèndix A: axiomàtica relacional.]
        Ofereix una presentació en què les fletxes són
        els individus primitius i els objectes es
        reconstrueixen a partir de les identitats.
        També hi trobareu la formalització del principi
        de dualitat, en versió sintàctica i semàntica.
        És una ampliació adequada després del primer
        capítol, sobretot per a qui s'interessi per
        la formalització lògica. Demana familiaritat
        amb lògica de primer ordre amb igualtat.
  \item[Apèndix B: grafs.]
        Aclareix les convencions, els homomorfismes,
        els camins i el pas a categories. Consulteu-lo
        quan treballar amb camins, llaços o arestes
        múltiples us plantegi dubtes en els primers
        capítols.
  \item[Apèndix C: convenció de mida.]
        Explica el marc de conjunts i classes del llibre.
        És especialment útil quan apareixen
        $\cat{Cat}$, categories de functors i
        categories de prefeixos. Cal distingir
        \enquote{petita} de \enquote{localment petita}: la segona
        condició garanteix homconjunts, no que tots
        els objectes formin un conjunt.
  \item[Apèndix D: solucions dels tests.]
        Respostes i justificacions. Consulteu-les només
        després d'haver respost tot el test.
\ifpublicacioparcial\else
  \item[Apèndix E: quasiexemples i contraexemples.]
        Recull estructures que fallen exactament un axioma
        (categoria, functor, transformació natural), la
        jerarquia entre igualtat, isomorfisme, isomorfisme
        natural i equivalència, i models de la lògica que
        queden fora dels topos. La secció sobre la
        jerarquia és útil des del capítol~\ref{cap:transformacions};
        la de la lògica, després del capítol~\ref{cap:topos}.
\fi
\end{description}

La taula de notació serveix per recuperar el significat
i la primera aparició dels símbols. El glossari
català--anglès ajuda a llegir la bibliografia internacional,
sobretot en termes com \emph{pullback}, \emph{slice},
\emph{presheaf} i \emph{whiskering}. L'índex alfabètic
és una eina per retrobar conceptes quan reapareixen en
un context nou.

Les lectures recomanades al final dels capítols tenen
una funció de contrast i ampliació. En una primera
lectura, és millor consultar una font alternativa per
aclarir una pregunta concreta que obrir alhora molts
tractaments generals del mateix tema.

\section{Una manera de mesurar l'aprenentatge}

No mesureu el progrés només pel nombre de pàgines llegides.
Una fita més fiable és poder fer, sense el text davant,
quatre operacions amb una idea nova:
\begin{enumerate}
  \item enunciar-ne la definició amb les hipòtesis correctes;
  \item interpretar-la en almenys dos exemples;
  \item utilitzar-la en un argument o un càlcul senzill;
  \item explicar què falla en un contraexemple o quan
        es retira una hipòtesi.
\end{enumerate}

No cal satisfer aquestes quatre condicions amb la mateixa
fluïdesa des del primer dia. És raonable avançar amb
una dificultat ben identificada i tornar-hi després.
El que convé evitar és substituir la comprensió per
la familiaritat visual amb els símbols.

En acabar el llibre, el guany més important serà poder
relacionar les peces: veure una propietat universal
darrere d'una construcció, una naturalitat darrere
d'una compatibilitat, una adjunció darrere d'una
correspondència i una operació sobre subobjectes
darrere d'una lectura lògica. Aquesta capacitat de
passar entre exemples i estructura és el criteri
que dona unitat a tot el recorregut.

\endgroup
